{\large\textbf{The Stern-Gerlach Experiment}}\footnote{H. S. Mani (former Director, HRI, Prayagraj) and Gautam Datta (DAIICT, Gandhinagar) were the principal authors of this problem. The contributions of the Academic Committee, Academic Development Group, and the International Board are gratefully acknowledged.}

\textbf{The Stern-Gerlach experiment} was performed in 1922 and it eventually
led to the determination of the magnetic moment of the electron. In this
experiment a beam of silver atoms each of mass
$m = 1.80 \times 10^{-25}$ kg emerges from an oven kept at temperature
T=1.20 $\times$ 10$^{3}$ K (see Figure 1). Assume that on emerging from the
oven all the atoms have the same momentum along the direction of the beam
($z$-direction). Ignore gravity.

\begin{center}\includegraphics[width=0.95\linewidth]{figures/APhO_2022_Q1_fig1.png}\end{center}

\begin{center}Figure 1: Schematic diagram of the Stern-Gerlach setup.\end{center}

\textbf{A.1}\quad \textbf{Speed of the Silver Atoms:} The speed $v_z$ of the
silver atoms emerging from the oven can be estimated to be $\sqrt{3k_BT/m}$
using the equipartition theorem. Calculate this value. \hfill 0.5pt

\textbf{B.1}\quad \textbf{The Basic Expression:} After emerging from the oven
the silver atoms move along the $z$- direction over a distance
$l_1 = 0.25$ m. Next, the silver atoms pass between two magnets over a
distance $l_2 = 0.5$ m. The magnets produce inhomogeneous magnetic field $B$
in the $x$-direction with constant gradient $dB/dx$. Assume that the silver
atom has a magnetic moment pointing either in + $x$-direction or -
$x$-direction, i.e. $\vec{\mu}_s = \pm\mu_s\hat{\imath}$. After passing
through the magnets the silver atoms pass through a further distance
$l_3 = 0.25$ m before striking the screen $PP'$. The distance between the two
striking beams on the screen is $\Delta x$. Derive the expression for the
splitting distance $\Delta x$ on the screen. \hfill 2pt

\textbf{The Inhomogeneous Magnetic Field:}

This part is concerned with the setup to create the inhomogenous magnetic
field ($dB/dx \neq 0$). It consists of a number of sub-parts. Two very long
wires parallel to the $z$-axis carry currents of magnitude $I_0$ and are
located at $A_1\,(0, -a, z)$ and $A_2\,(0, a, z)$ (see figure below). The
direction of the current passing through $y = -a$ is $-\hat{k}$ and for the
one passing through $y = a$ is $\hat{k}$. The entire system is inside a medium
of high relative magnetic permeability $\mu_r$. We take $\mu$ = $\mu_0\mu_r$.
The wires are insulated and no current leaks into the medium.

\begin{center}\includegraphics[width=0.45\linewidth]{figures/APhO_2022_Q1_fig2.png}\end{center}

\begin{center}Figure 2: The arrangement for the inhomogeneous magnetic field\end{center}

\textbf{C.1}\quad Derive the expression for the magnetic field vector at a
point $P_1\,(x, y, 0)$ in the $x$-$y$ plane (see Fig. 2). \hfill 1.5pt

\textbf{C.2}\quad Consider a circle with its centre C on the $x$-axis
$(x_c, 0)$ and with radius equal to the distance $AC$ (see Figure 2). Obtain
the direction of the magnetic field at circumferential points $R$ (on the $x$
axis) and at $P_0$ ($CP_0$ is parallel to the $y$ axis). \hfill 0.5pt

\textbf{C.3}\quad Now a small slice of the high magnetic permeability medium
between the circles with radii AC and AD is removed and replaced with air at
very low pressure (see Fig. 2). One can show from continuity considerations
that the magnetic field in this gap region is given by the same expression as
if the magnetic medium were not removed. (We shall assume this and you are
not required to prove this). Hence state the expression for the magnetic field
at the point $(x, 0)$ in the air gap region. \hfill 0.5pt

\textbf{D.1}\quad \textbf{The Force:} As mentioned earlier the silver atoms are
travelling in the $(x, 0, z)$ plane with their velocities parallel to the
z-axis and given by $\vec{v} = v_z\hat{k}$. Recall also that the magnetic
dipole of the silver atom is $\vec{\mu}_s = \pm\mu_s\hat{\imath}$. Obtain the
expression for the magnitude of the force $F_x$ acting on a silver atom along
the x-direction in terms of $\mu_s$ , $I_0$ , $a$ $\mu$ and relevant
coordinates. \hfill 0.5pt

\textbf{E.1}\quad \textbf{The Magnetic Field and its Gradient:} We assume that
this same force $F_x$ acts over a small distance $l_2$ along the $z$ axis
(Figure 1). Assume also that the silver atoms pass through the mid-point $P$
of $RQ$ (Figure 2). The following experimental values are given:
\[
\frac{\mu}{\mu_0} = 10^4\,;\quad a = 0.60\,cm\,;\quad OC = 0.60cm\,;\quad
OD = 0.80\,cm\,; I_0 = 2.00A
\]
Here $\mu_o$ is the magnetic permeability of free space. Obtain the numerical
value of the magnitudes of the magnetic field $B_P$ and its gradient
$dB_P/dx$ at this mid-point in S.I. units. \hfill 2.0pt

\textbf{F.1}\quad \textbf{The Magnetic Moment of the Silver Atom:} If
$v_z = 500\,\mathrm{m{\cdot}s^{-1}}$ and magnetic field is calculates as
above, the Stern-Gerlach experiment yields a split of $\Delta x = 0.20$ cm.
Obtain the value of the magnetic moment of the silver atom $\mu_s$ in S.I.
units. \hfill 1.5pt

\textbf{G.1}\quad \textbf{The spread of the line:} The silver atoms may not
all have the same speed. Let there be a spread of 20 \% in the beam speed.
What would be the consequent spread of the dot $\delta x$ on the screen?
\hfill 0.5pt

\textbf{H.1}\quad \textbf{The error in the magnetic moment:} What is the
consequent error bar on the evaluation of the magnetic moment $\delta\mu_s$?
\hfill 0.5pt
